\begingroup
\ArithmeticHead{248}{L}{ANALYSIS. BRIEFWECHSEL.}
\SourceLeaf{255b}
\pdfbookmark[1]{[8.] Gauss an Bessel, 30. März 1828, mit Anhang}{x1-correspondence-8}
\NumberedItem{8}
\begin{center}\textsc{Gauss} an \textsc{Bessel}. Göttingen den 30. März 1828.\end{center}
\HistoricCitation{[Briefwechsel zwischen \textsc{Gauss} und \textsc{Bessel}, Leipzig 1880, S. 477.]}
\noindent .\enspace .\enspace .\enspace .\enspace .\enspace .\quad Zur Ausarbeitung der seit vielen Jahren (1798) angestellten Untersuchungen über die transcendenten Functionen werde ich vorerst wol noch nicht kommen können, da erst noch mit manchen andern Dingen aufgeräumt werden muss. Hr. \textsc{Abel} ist mir, wie ich sehe, jetzt zuvorgekommen\textsuperscript{[*)]} und überhebt mich in Beziehung auf etwa ${}^1\! /\! {}_3$ dieser Sachen der Mühe, zumahl da er alle Entwickelungen mit Eleganz und Concision gemacht hat. Er hat gerade denselben Weg genommen, welchen ich 1798 einschlug, daher die grosse Übereinstimmung der Resultate nicht zu verwundern ist. Zu meiner Bewunderung erstreckt sich dies sogar auf die Form und zum Theil auf die Wahl der Zeichen, so dass manche seiner Formeln wie eine \textls[120]{reine Abschrift} der meinigen erscheinen. Jeder Misdeutung zuvorzukommen bemerke ich jedoch, dass ich mich nicht erinnere, von diesen Sachen irgend jemanden etwas mitgetheilt zu haben.

\par\indent Nicht dasselbe kann ich aber von einer Kleinigkeit sagen, die Sie im 1. Heft des 3. Bds von \textsc{Crelle}’s Journal\textsuperscript{[**)]} von mir finden werden.
\par\vfill\noindent\rule{.25\linewidth}{.4pt}\par\smallskip
\begingroup\footnotesize
\noindent [*) Siehe N. H. \textsc{Abel}, \textit{Recherches sur les fonctions elliptiques}, \textsc{Crelles} Journal für Mathematik \textbf{2} (1828), S. 101, Oeuvres, nouvelle édition, I, 1881, S. 263—351; dieser erste Teil der \textit{Recherches} ist am 20. September 1827 erschienen.]\par
\noindent [**) \textit{Beweis eines algebraischen Lehrsatzes}, \textsc{Crelles} Journal für Mathematik \textbf{3}, 1828, S. 1, Werke III, S. 65.]\par\endgroup
\clearpage
\ArithmeticHead{249}{R}{GAUSS AN BESSEL.}
\SourceLeaf{256}
\noindent Diesen Beweis des \textsc{Harriot}’schen Lehrsatz[es] habe ich, seitdem ich in Göttingen bin, sehr häufig mündlich vorgetragen\textsuperscript{[*)]}; der verstorbene \textsc{Pfaff} schrieb mir u. a. mehrere Jahre vor seinem Tode\textsuperscript{[**)]}, dass er ihm in Halle von jemand (als von mir herrührend) mitgetheilt sei, der ihn wieder durch die zweite oder dritte Hand aus Jena erhalten habe, wie ich vermuthe von dem verst[orbenen] \textsc{Posselt}; \textsc{Pfaff}, dessen Beifall dieser Beweis erhalten hatte, ersuchte mich um meine Zustimmung, um ihn in einer Sammlung von Abhandlungen, die er herauszugeben beabsichtigte, bekannt zu machen. Ich gab ihm diese natürlich gern; allein das projectirte Buch ist meines Wissens nicht erschienen. Ob dies in einem Zusammenhange mit dem Erscheinen eines Beweises von derselben Grundidee in dem letzten Hefte des $2^{\mathrm{n}}$ Bandes von \textsc{Crelles} Zeitschrift\textsuperscript{[***)]} (welches mir erst fast zugleich mit dem $1^{\mathrm{ten}}$ des 3. Bdes zu gesichte gekommen ist, etwa 3 oder 4 Monate später als ich \textsc{Crelle} den Aufsatz geschickt hatte) steht, muss ich auf sich beruhen lassen.\quad .\enspace .\enspace .\enspace .\enspace .\enspace .

\begin{center}\large Anhang zu [8.]\end{center}
\begin{center}[a.]\par\medskip \textsc{Pfaff} an \textsc{Gauss}. Halle, 20. Oktober 1824.\end{center}
\noindent \{.\enspace .\enspace .\enspace .\enspace .\quad Ehe ich meinen alten Vorsatz, ein Lehrbuch der algebraischen und höheren Analysis (über welche ich öfter lese) herauszugeben, ausführe, gedenke ich einige Abhandlungen und weitere Ausführungen einzelner Lehren (wie auch \textsc{Kaestner} gerathen und gethan hat) herauszugeben, und vielleicht schon in der nächsten Ostermesse einen Anfang damit zu machen. Darunter gehört auch die Entwickelung des \textsc{Harriot}schen Satzes. Zufällig erhielt ich vor einigen Jahren von einem Zuhörer, der früher in Jena Analysis bey Prof. v. \textsc{Münchow} gehört hatte, ein Vorlesungs-Heft von letzterem, worin auch ein Beweis des \textsc{Harriot}schen Satzes, als von Ihnen herrührend und mitgetheilt, vorkommt. In ein schon seit geraumer Zeit ins reine gearbeitetes Manuscript
\par\vfill\noindent\rule{.25\linewidth}{.4pt}\par\smallskip
\begingroup\footnotesize
\noindent [*) Vergl. M. A. \textsc{Stern}, \textit{Über einen Satz von Gauss}, Göttinger Nachrichten 1869, S. 330, abgedruckt Werke X2, S. 70.]\par
\noindent [**) Siehe die als \textls[120]{Anhang} zu diesem Briefe abgedruckten beiden Briefstellen.]\par
\noindent [***) Gemeint ist jedenfalls: \textit{Beweis des Harriotschen Satzes} von J. A. \textsc{Grunert}, \textsc{Crelles} Journal für Mathematik \textbf{2} (1828), S. 335—344.]\par\endgroup
\smallskip\noindent\small X1.\hfill 32\normalsize
\clearpage
\ArithmeticHead{250}{L}{ANALYSIS. BRIEFWECHSEL.}
\SourceLeaf{257}
\noindent nahm ich (mit denjenigen Änderungen, welche die Verbindung mit dem übrigen erforderte) Ihren einfachen und sinnreichen Beweis auf, weil durch dessen Weglassung eine wesentliche Lücke in der Abhandlung entstanden wäre, und weil ich dachte, dass Sie bei Ihren übrigen wichtigen und schweren Untersuchungen schwerlich so bald dazu kommen würden, diesen Beweis selbst bekannt zu machen. Jetzt benuze ich gegenwärtige Gelegenheit, Sie um die Genehmigung dieser Bekanntmachung zu bitten.\quad .\enspace .\enspace .\enspace .\enspace .\enspace .\enspace .\enspace .\enspace .\enspace .\enspace .\}

\begin{center}[b.]\par\medskip \textsc{Gauss} an \textsc{Pfaff}. Göttingen, 21. März 1825.\end{center}
\HistoricCitation{[Sammlung von Briefen gewechselt zwischen \textsc{Joh. Friedrich Pfaff} und .\enspace .\enspace .\enspace .\ Anderen; herausg. von Dr. \textsc{Carl Pfaff}, Leipzig 1853, S. 277—278.]}
\noindent .\enspace .\enspace .\enspace .\enspace .\enspace .\quad Im nächsten Sommer werden mich meine Dreiecksmessungen wieder von Göttingen entfernen; ich habe noch ein Dreiecksnetz von Barmen bis Ostfriesland, der Nordsee und Helgoland zu führen. Ich wünsche sehr alle Arbeiten dieser Art, die noch rückständig sind, in einem Stück zu vollenden, um dann die Lebensjahre, die der Himmel mir noch schenken wird, ungestört auf Arbeiten im Studirzimmer wenden zu können. Jene Messungen erhalten einen ihrer vornehmsten Reize für mich durch den Umstand, dass sie mir Gelegenheit zur Anwendung der mir eigenthümlichen theoretischen Methode, solche Messungen zu behandeln, geben. Nach Beendigung der Messungen werde ich darüber ein eigenes Werk, vermuthlich von bedeutender Ausdehnung ausarbeiten\textsuperscript{[*)]}. Meine Preisschrift über die Transformation der Flächen ist, wie ich höre, jetzt unter der Presse\textsuperscript{[**)]}. Gegenwärtig bin ich mit einer Untersuchung aus der höheren Arithmetik beschäftigt als Anfang der Theorie der biquadratischen Reste; wenn es möglich ist, werde ich die Abhandlung der
\par\vfill\noindent\rule{.25\linewidth}{.4pt}\par\smallskip
\begingroup\footnotesize
\noindent [*) Erschienen sind die Abhandlungen\par
\noindent\hspace{1.5em}\textit{Disquisitiones generales circa superficies curvas}, 1827, Werke IV, S. 217;\par
\noindent\hspace{1.5em}\textit{Untersuchungen über Gegenstände der höheren Geodaesie} I, 1843, ebenda, S. 259; II, 1846, ebenda, S. 301.]\par
\noindent [**) \textit{Allgemeine Auflösung der Aufgabe, die Theile einer gegebenen Fläche auf einer andern gegebenen Fläche so abzubilden, dass die Abbildung dem abgebildeten in den kleinsten Theilen ähnlich wird}, erschienen Altona 1825, Werke IV, S. 189; vergl. in bezug auf den Druck dieser Abhandlung \textit{Briefwechsel zwischen C. F. Gauss und H. Schumacher} II, 1860. S. 5—7.]\par\endgroup
\clearpage
\ArithmeticHead{251}{R}{GAUSS AN PFAFF.}
\SourceLeaf{258}
\noindent k. Societät noch vor meiner diesjährigen Abreise vorlegen\textsuperscript{[*)]}. Mit Verlangen sehe ich Ihrem neuen Werk entgegen\textsuperscript{[**)]}. Ich rechne mir’s zur Ehre, dass Sie meinen Beweis des \textsc{Harriot}schen Lehrsatzes darin aufnehmen wollen: ich habe ihn öfters mündlich vorgetragen, ohne doch zu wissen, wie er in die Hände des Professor \textsc{Münchow} gekommen ist. Die Besorgniss, dass er, durch mehrere Hände gegangen, vielleicht verunstaltet sein könnte, beunruhigt mich nicht, da Sie in diesem Fall dies sogleich erkannt haben und mir ihn dann erst zur Anerkennung communicirt haben würden. So wie ich mir ihn aufgezeichnet habe, füllt er nicht viel über eine Octavseite.\quad .\enspace .\enspace .\enspace .\enspace .\enspace .\enspace .
\par\vfill\noindent\rule{.25\linewidth}{.4pt}\par\smallskip
\begingroup\footnotesize
\noindent [*) Die Abhandlung \textit{Theoria residuorum biquadraticorum, Comment. prima}, Werke II, S. 65, ist der Königl. Societät am 5. April 1825 vorgelegt worden. Dem Briefwechsel nach hat \textsc{Gauss} seine Reise um den 20. April angetreten.]\par
\noindent [**) Dieses Werk von \textsc{Pfaff} ist, vergl. \textsc{Gauss} an \textsc{Bessel}, oben S. 247, nicht erschienen; \textsc{Pfaff} ist am 20. April 1825 gestorben.]\par\endgroup
\clearpage
\endgroup
